\begin{figure*}
  \includegraphics[width=\textwidth]{img/eye_catch_ACMMM.pdf}
  \caption{
    Overview of the Script2Layout task. From left to right: (1) a manga page from the Manga109 dataset; (2) its corresponding script, generated using a vision large language model (VLLM) with visual prompting and speaker-aware dialogue sequencing; and (3) the layout predicted by a Script2Layout model based on the generated script. The script follows standard manga production conventions and includes scene headings, action descriptions, and dialogue. This setup allows layout models to make use of narrative structure for generating more coherent and expressive layouts. (Source: *Hanzai Kousyounin Minegishi Eitarou* by Takashi Kii)
  }
  \label{fig:eyecatch}
  \vspace{-1.2em}
\end{figure*}

\section{Introduction}
Storytelling takes many forms, from text-focused novels to visually-driven animation. Among these, manga has developed a unique expressive form that divides a flat surface into panels containing characters and dialogue. In typical manga production, the process of translating a narrative into a visual layout occurs in a preliminary draft called a \emph{name}. The layout therefore embodies critical directorial decisions. Its composition, particularly the positioning of panels, characters, and speech balloons, is a deliberate choice designed to structure the visual narrative and ensure the story is communicated clearly and with intended emotional weight~\cite{fukaya_mangalayout}.

In contrast to this narrative-centric process, the field of layout generation has seen development, which has primarily focused on functional applications like Web UIs~\cite{Kikuchi2021_webUI_layout} and posters~\cite{tanaka_scipostlayout_2024, advertising_posters}, aiming to generate visually compelling layouts that ensure clear and efficient information transfer~\cite{shi_intelligent_2023}. However, layouts that combine storytelling and dramatic effects, like manga, are under-researched. The absence of data linking stories to layouts is a key reason for this gap. Script data are often produced internally and rarely released publicly, making it difficult to obtain. This data scarcity has been the primary barrier to developing assistive applications for this narrative-driven creative domain.

In this study, we utilized Vision Large Language Models (VLLMs) to create Manga109Script, a script dataset for approximately 20,000 manga pages based on the Manga109 dataset~\cite{matsui2017sketch,aizawa_building_2020}.
Recent advances in VLLMs have enabled high-precision extraction of textual information from images~\cite{zhang_vision-language_2024}. This capability has significantly contributed to research in Visual Storytelling, which extracts narratives from image sequences~\cite{oliveira_story_2024, li_benchmark_2025}.
Despite advancements in VLLMs, understanding content across multiple panels remains challenging~\cite{ikuta_mangaub_2024,vivoli_comicap_2024}.
To mitigate VLLMs' limitations, we applied Visual Prompting~\cite{shtedritski_what_2023} with the Manga109 dataset annotations and employed a rule-based algorithm designed for panel reading order~\cite{sachdeva_manga_2024}, to generate high-quality narrative script data.

Our scripts follow standard script formatting conventions, including {\it scene headings} indicating location, {\it action descriptions} detailing character movements and scenery, and {\it dialogue} (Fig.~\ref{fig:eyecatch}).
We consulted with manga experts to confirm the format fits professional workflows.
The generated script's quality was also evaluated by humans.

To demonstrate the practical utility of our dataset for such creative applications, we developed manga layout generation models that take scripts as direct input.
We comprehensively evaluated these models against existing layout generation methods and assessed the impact of script input to establish a baseline for the Script2Layout task. Dataset, and code will be released upon publication.

Our contribution is threefold:
\begin{enumerate}
  \item We introduce an annotation-assisted method to generate script data from manga images using VLLMs, validated by human evaluators.
  \item We publish the first professional script-format manga storytelling dataset, covering 20,000 pages.
  \item We establish robust baselines for the novel Script2Layout task, demonstrating for the first time how script influences manga layout.
\end{enumerate}

